% chapters/13/section-03.tex
\section{Kontrol Pemain}
\label{sec:kontrol-pemain}

\subsection{Kode Lengkap}

\kodecsharp{unity/ch13_kontrol_pemain.cs}{Kontrol pemain platformer 2D berbasis Rigidbody2D}{lst:ch13-pemain}

\subsection{Mengapa Dua Fungsi, bukan Satu}

Skrip itu memuat dua fungsi: \code{Update} dan
\code{FixedUpdate}. Pembagian pekerjaan di antara keduanya bukan pilihan gaya
penulisan, melainkan keharusan yang lahir dari sifat keduanya.

\begin{itemize}
  \item \textbf{\code{Update} dijalankan sekali setiap bingkai}, dengan selang
        waktu yang berubah-ubah. Fungsi ini dipakai untuk membaca masukan,
        sebab masukan pemain berkaitan dengan bingkai yang dilihatnya, dan
        dapat berubah beberapa kali dalam satu bingkai fisika.
  \item \textbf{\code{FixedUpdate} dijalankan dengan selang waktu tetap},
        mengikuti pengaturan \istilah{fixed timestep} pada mesin. Fungsi ini
        dipakai untuk menggerakkan \istilah{Rigidbody} dan untuk memeriksa
        tanah, sebab keduanya berkaitan dengan perhitungan fisika.
\end{itemize}

Pembagian itu persis sama dengan pembedaan langkah berubah dan langkah tetap
yang telah Anda pelajari pada Bagian~2.2. Yang berbeda hanyalah cara
mewujudkannya: pada Pygame Anda menuliskannya sendiri, sedangkan pada Unity
mesin menyediakan dua fungsi yang terpisah.

\begin{peringatanbox}
Memindahkan \istilah{Rigidbody} melalui \code{transform.position} pada
\code{Update}---sebagaimana biasa dilakukan pada blok Python---akan
mengakibatkan benda menembus dinding dan bergetar. Mesin menghitung fisika pada
\code{FixedUpdate}, sehingga perubahan letak yang dilakukan di luarnya tidak
diperhitungkan. Gerakkanlah \istilah{Rigidbody} melalui komponennya, dan
kerjakanlah pada \code{FixedUpdate}.
\end{peringatanbox}

\subsection{Pemeriksaan Tanah}

Perhatikan cara memeriksa apakah pemain sedang menyentuh tanah pada skrip di
atas. Pemeriksaannya memakai \code{Physics2D.OverlapCircle}, yaitu memeriksa
apakah ada benda lain pada sebuah lingkaran kecil di bawah kaki pemain.

Cara itu dipilih alih-alih membandingkan posisi dengan letak lantai, sebab
lantainya tidak harus berupa satu bidang datar. Lingkaran kecil itu bekerja
sama baiknya pada lantai mendatar, lantai bertingkat, maupun bongkahan miring.

Perhatikan pula bahwa lingkaran itu digambar pada tampilan \istilah{Scene}
melalui fungsi \code{OnDrawGizmosSelected}. Kebiasaan menggambar daerah
pemeriksaan semacam itu sangat berguna: tanpa gambarnya, sulit diketahui apakah
lingkaran pemeriksaannya terlalu kecil, terlalu besar, atau salah letak.

\subsection{Membaca Masukan dan Gerak}

Perhatikan tiga hal pada cara masukan dan gerak diperlakukan.

\begin{enumerate}
  \item \textbf{Masukan dibaca pada \code{Update} dan disimpan pada peubah},
        lalu dipakai pada \code{FixedUpdate}. Cara itu mencegah masukan yang
        terlewat, yaitu penekanan tombol yang terjadi di antara dua langkah
        fisika.
  \item \textbf{Lompatan dibaca sebagai peristiwa} melalui
        \code{Input.GetButtonDown}, bukan sebagai keadaan tombol. Alasannya
        sama dengan yang dijelaskan pada Bagian~6.3: menahan tombol tidak boleh
        menghasilkan lompatan berulang.
  \item \textbf{Kecepatan tegak dipertahankan ketika gerak mendatar diubah.}
        Baris gerak mendatar hanya mengganti komponen
        mendatarnya, sedangkan komponen tegaknya dibiarkan. Menggantinya
        seluruhnya akan menghapus pengaruh gravitasi.
\end{enumerate}

\begin{tipsbox}
Perhatikan baris \code{[RequireComponent(typeof(Rigidbody2D))]} pada awal
skrip. Baris itu memberitahu Unity bahwa skrip tersebut menuntut keberadaan
\istilah{Rigidbody2D}, sehingga komponennya ditambahkan secara otomatis
apabila belum ada. Kebiasaan semacam itu mencegah kesalahan yang paling sering
terjadi pada tahap ini, yaitu skrip yang menempel pada benda tanpa
\istilah{Rigidbody}.
\end{tipsbox}
